\documentclass[10pt,a4paper]{amsart}
\usepackage[margin=19mm]{geometry}
\usepackage{amsmath,amssymb,mathtools}
\usepackage[scr=rsfs,cal=euler]{mathalfa}
\usepackage{xfrac}
\usepackage[unicode,hidelinks,bookmarks=false]{hyperref}
\newcommand{\ie}{i.e.\ }
\newcommand{\cf}{cf.\ }
\newcommand{\resp}{respectively\ }
\newcommand{\Subsection}{\subsection}
\providecommand{\nobreakdash}{\nobreak\hbox{-}}
\setlength{\emergencystretch}{2em}
\multlinegap0pt
\newcommand{\C}{\mathbb C}
\newcommand{\HH}{\mathbb H}
\newcommand{\N}{\mathbb N}
\newcommand{\Z}{\mathbb Z}
\newcommand{\calT}{\mathcal T}
\newcommand{\dd}{\delta}
\newcommand{\thetac}[2]{{\vartheta\left[\begin{smallmatrix}
        #1 \\ #2
    \end{smallmatrix}\right]}}
\theoremstyle{plain}
\newtheorem{definition}{Definition}[section]
\newtheorem{theorem}[definition]{Theorem}
\newtheorem{lemma}[definition]{Lemma}
\newtheorem{corollary}[definition]{Corollary}
\newtheorem{proposition}[definition]{Proposition}
\newtheorem{remark}[definition]{Remark}
\newtheorem*{theorem*}{Theorem}
\title[Derivatives of Theta Functions and a Problem of Lyubarskii and Nes]{Derivatives of Theta Functions and a Problem of Lyubarskii and Nes}
\author{Lukas Liehr, Irina Shafkulovska, Mitchell A. Taylor}
\date{Source: arXiv:2609.27846v1 (CC BY 4.0)}
\begin{document}
\maketitle

\noindent\emph{Source and licence.} The delivered work is the article shown in the title above, version arXiv:2609.27846v1 (CC BY 4.0), distributed under the Creative Commons Attribution 4.0 International licence, as recorded in the arXiv licence metadata for this exact version and shown on the abstract page saved with this item. The full licence notice, which carries the licence URL and the address of that abstract page, is the bundled file \texttt{licenses/arxiv-2609.27846-CC-BY-4.0.txt}.
\par\smallskip

\noindent\textit{Editorial scope.} Counted unit: the article's Theorem (source label thm:primitive-derivative) (source0.tex lines 441--456, source label \texttt{thm:primitive-derivative}): ``Let and satisfy (p,q) 1, q n p 2. For every , every and every polynomial of degree , the map p,a,b ( ) q, f ( P( z) f (w r q ) ) r 0 q-1 is injective.''. It is delivered together with the article's complete original proof (source0.tex lines 458--622, one proof environment adjacent to the statement, 165 lines), copied line for line from the article's own text. Required context retained uncounted: eq:thetac\_quasip (lines 287--291); lem:torsion-jet (lines 360--374); lem:frobenius (lines 398--415). Editorial formatting only: an AMS \texttt{amsart} preamble that restates the article's own macro definitions and its own theorem declarations, and the theorem environments the retained lines use; the article's own class file is neither shipped nor input; the retained environments are renumbered sequentially in order of appearance, whereas the article prints them with its own numbering; the article's equation numbering is not reproduced and every cross-reference inside the retained text resolves to the wrapper's numbering. No citation occurs inside any retained line, so the wrapper carries no bibliography; All retained bodies are exact copies of the bundled \texttt{sources/211520/source0.tex}, so no omission receipt applies. No asset-inclusion command occurs inside any retained span and the package ships no image asset.


\section*{Retained context: eq:thetac\_quasip (source0.tex lines 287--291)}

\begin{equation}\label{eq:thetac_quasip}
\thetac{a}{b}(z+1,\tau)=e^{2\pi i a}\thetac{a}{b}(z,\tau),
\ \
\thetac{a}{b}(z+\tau,\tau)=e^{-\pi i(2b+2z +\tau)}\thetac{a}{b}(z,\tau).
\end{equation}


\section*{Retained context: lem:torsion-jet (source0.tex lines 360--374)}

\begin{lemma}\label{lem:torsion-jet}
Let $\tau\in\HH$, $q\in\N, \, q\geq 2$ and let 
$\Theta(z) = \thetac{a}{b}(z,\tau)$ for some $a,b\in\C$.
Let $P$ be a
nonzero polynomial with
$$
\deg P\leq q-2.
$$
Then, for every $w\in\C$, the system of equations
\begin{equation}\label{eq:impossible_system}
    P(\dd_z)\Theta\left(w+\frac{j}{q}\right)=0,
\quad j=0,\dots,q-1,
\end{equation}
cannot be satisfied.
\end{lemma}


\section*{Retained context: lem:frobenius (source0.tex lines 398--415)}

\begin{lemma}
\label{lem:frobenius}
Let $p\in\N$, $b\in\C$ and $\tau\in\HH$.
Let $\phi_0,\dots,\phi_{p-1}$ be a basis of $\calT_{p,0,b}(\tau)$. Define
$$
\Delta(z_0,\dots,z_{p-1})
:=
\det[\phi_s(z_l)]_{l,s=0}^{p-1},
\quad
S(z):=\sum_{l=0}^{p-1}z_{l}.
$$
Then there exists a nonzero constant $C$ such that
$$
\Delta(z_0,\dots,z_{p-1})
=
C\, \thetac{1/2+p/2}{1/2+b+p/2}(S(z), \tau) \prod\limits_{0\leq l <l'< p}\thetac{1/2}{1/2}(z_{l'}-z_l,\tau).
$$
\end{lemma}


\section{The counted unit: Theorem (source label thm:primitive-derivative) and its complete original proof}

\begin{theorem}\label{thm:primitive-derivative}
Let $a,b\in\C$ and $n,p,q\in\N$ satisfy
$$
(p,q)=1,
\quad
q\geq {n}p+2.
$$
For every $w\in\C$, every $\tau \in \mathbb H$ and every polynomial $P$ of degree $n$, the map
$$
\calT_{p,a,b}(\tau)\to \C^q,
\quad
f \mapsto
\left({P(\dd_z)} f\left(w+\tfrac{r}{q}\right)\right)_{r=0}^{q-1}
$$
is injective.
\end{theorem}

\begin{proof}
{
Since 
\begin{equation}
    \calT_{p,a,b}\to \calT_{p,0,b+a\tau}, \quad f\mapsto e^{-2\pi i a\cdot} f
\end{equation}
is an isomorphism and 
\begin{equation}
    \dd_z f(z) = \dd_z \left(e^{2\pi i az}g(z)\right)= e^{2\pi i a z}(a g(z) + \dd_z g(z)),\quad g(z) = e^{-2\pi i az} f(z),
\end{equation}
by Leibniz' rule there exists a polynomial $Q$ of degree at most $n$ such that
\begin{equation}\label{eq:P_vs_Q_for_a}
    P(\dd_z) f (z) = e^{2\pi i a z} Q(\dd_z) g(z),\quad g(z) = e^{-2\pi i az} f(z). 
\end{equation}
The polynomial $Q$ has degree $n$ because its leading coefficient equals the leading coefficient of $P$. Thus, we can assume without loss of generality that $a=0$ and that $P$ has a leading coefficient equal to one.
}

{Let $\phi_0,\dots,\phi_{p-1}$ be a basis of $\calT_{p,0,b}(\tau)$.}  Suppose, to
the contrary, that the stated map is not injective.  Then the
$q\times p$ matrix
$$
R_{r,s}:={{P(\dd_z)}}\phi_s\left(w+\tfrac{r}{q}\right),
\quad
r=0,\dots,q-1,
\quad
s=0,\dots,p-1,
$$
has rank strictly smaller than $p$.  Hence, all of its $p\times p$ minors
vanish.

For $j\in\Z$, define the consecutive minor
$$
M_j:=
\det\left[
{P(\dd_z)}\phi_s\left(w+\tfrac{j+l}{q}\right)
\right]_{l,s=0}^{p-1}.
$$
Each $P(\dd_z)\phi_s$ is $1$-periodic, so reducing the indices
$j+l$ modulo $q$ only permutes the rows. 
Rank failure gives
$$
M_j=0,
\quad j=0,\dots,q-1.
$$
Let
$$
\Delta(z_0,\dots,z_{p-1})
:=
\det[\phi_s(z_l)]_{l,s=0}^{p-1}.
$$
By Lemma~\ref{lem:frobenius}, we have
$$
\Delta(z_0,\dots,z_{p-1})
=
C\, \thetac{1/2+p/2}{1/2+b+p/2}(S(z), \tau)V(z_0,\dots,z_{p-1}) ,
$$
where
$$
S(z_0,\dots,z_{p-1}):=\sum_{l=0}^{p-1} z_l,
\quad
V(z_0,\dots,z_{p-1})
:=\prod\limits_{0\leq l<l'< p}\thetac{1/2}{1/2}(z_{l'}-z_l,\tau).
$$
We apply $P(\dd_z)$ to each row of $[\phi_s(z_l)]_{l,s=0}^{p-1}$ with respect to $z_l$ and denote this by 
$$
P(\dd_{z_0})\cdots P(\dd_{z_{p-1}}) [\phi_s(z_l)]_{l,s=0}^{p-1}.
$$
We then restrict the argument to 
$$
z_l=w+\tfrac{l}{q},
\quad l=0,\dots,p-1.
$$
On the left, multilinearity of the determinant gives
$$
\left.
P(\dd_{z_0})\cdots P(\dd_{z_{p-1}})\Delta
\right|_{z_l=w+l/q}
=
\det\left[
P(\dd_z)\phi_s\left(w+\frac{l}{q}\right)
\right]_{l,s=0}^{p-1}
=:M(w).
$$
At the restricted points, we have
$$
z_{l'}-z_l=\tfrac{l'-l}{q}.
$$
As $1\leq l'-l\leq p-1<q$, none of these differences lies in
$\Z+\tau\Z$.  Since $\thetac{1/2}{1/2}(\cdot,\tau)$ has its only zero at the
origin modulo $\Z+\tau\Z$, it follows that
$$
V_0:=
V\left(w,w+\tfrac1q,\dots,w+\tfrac{p-1}{q}\right)
=
\prod\limits_{0\leq l <l'< p}\thetac{1/2}{1/2}(\tfrac{l'-l}{q},\tau)
\neq0.
$$
The above number is independent of $w$.
Let 
\begin{equation}
    \Theta(z)=\thetac{1/2+p/2}{1/2+b+p/2}(z,\tau).
\end{equation}
By Leibniz' rule, after the restriction $z_l=w+l/q$, every derivative of
$V$ is evaluated at the fixed tuple of differences
$\{(l'-l)/q\}_{l<l'}$ and is therefore a constant.  Hence, there are
constants $c_0,\dots,c_{np}$ such that
$$
M(w)
=
\sum_{d=0}^{np}c_d
\dd_z^d\Theta \left(pw+\tfrac{p(p-1)}{2q}\right).
$$
Equivalently,
$$
M(w)=Q(\dd_z)\Theta \left(pw+\tfrac{p(p-1)}{2q}\right),
$$
for a polynomial $Q$ of degree at most $np$.

The coefficient $c_{np}$ of the derivative of the highest order  comes only from the term in which the highest order
derivative part of every operator $P(\dd_{z_l})$
falls on the factor $\Theta(S(z_0,\dots,z_{p-1}))$ and no derivative falls on $V$. Therefore,
$$
c_{np}
=
C\prod_{0\leq l<l'< p}
\thetac{1/2}{1/2}\left(\tfrac{l'-l}{q},\tau\right)
=
CV_0
\neq0.
$$
Thus, $Q\neq0$ and $\deg Q={{n}}p$.

Since $M_j=M(w+j/q)$, the vanishing of the consecutive minors gives
$$
Q(\dd_z)\Theta\left(
pw+\tfrac{pj}{q}+\tfrac{p(p-1)}{2q}
\right)=0,
\quad j=0,\dots,q-1.
$$
Since $(p,q)=1$, for each $j$ there exists a residue $j'\in\{0,\dots,q-1\}$
and an integer $n_j$ such that
$$
\frac{pj}{q}=\frac{j'}{q}+n_j.
$$
Let
$$
w':=pw+\frac{p(p-1)}{2q}.
$$
The preceding vanishing can be expressed as
$$
Q(\dd_z) \Theta\left(w'+\tfrac{j'}{q}+n_j\right)=0,
\quad j'=0,\dots,q-1,
$$
after reordering the residues. 
By \eqref{eq:thetac_quasip}, $\Theta(z+n_j)=e^{2\pi i(1/2+p/2)n_j}\Theta(z)$ for all $z\in\C$, and hence $Q(\dd_z)\Theta(z+n_j)=e^{2\pi i(1/2+p/2)n_j}Q(\dd_z)\Theta(z)$. Since this factor is nonzero, it follows that
$$
Q(\dd_z)\Theta\left(w'+\tfrac{j}{q}\right)=0,
\quad j=0,\dots,q-1.
$$
Finally,
$$
\deg Q={{n}}p\leq q-2.
$$
Since $\Theta=\thetac{1/2+p/2}{1/2+b+p/2}(\cdot,\tau)$ and $q\geq np+2\geq 2$, this contradicts Lemma~\ref{lem:torsion-jet}. Hence, the original derivative evaluation map is injective.
\end{proof}

\end{document}
